\chapter{Python
编程基础}\label{ux9644ux5f55-3python-ux7f16ux7a0bux57faux7840}

\begin{quote}
\textbf{本章导读} 目标：用 Python 检查昇腾 310B
模型部署中常见的配置、清单和结果数据。
数据：正文内联示例，不依赖外部数据集。
标准库：\passthrough{\lstinline!json!}、\passthrough{\lstinline!statistics!}、\passthrough{\lstinline!math!}。
输出：可直接在 Python REPL 中运行，也可保存为
\passthrough{\lstinline!.py!} 文件复现。
\end{quote}

本附录按实际使用顺序讲解 Python。模型
ID、输入形状、预处理均值、精度模式和准入状态会反复出现在样例代码中，因此本章直接用这些对象编写示例。

\passthrough{\lstinline!for!}、\passthrough{\lstinline!if!}、\passthrough{\lstinline!dict!}、\passthrough{\lstinline!set!}、异常处理等基础语法会在前几节展开，第
4 节补充 \passthrough{\lstinline!samples/!}
中高频出现的文件、JSON、命令行和异步写法。所有示例都只使用标准库，不引入
\passthrough{\lstinline!pandas!}，也不需要板端运行环境。

\section{1.
数据、类型与容器}\label{ux6570ux636eux7c7bux578bux4e0eux5bb9ux5668}

\subsection{变量与基本类型}\label{ux53d8ux91cfux4e0eux57faux672cux7c7bux578b}

变量是给值起的名字。Python
中每个值都有类型，类型决定它可以参与哪些运算。检查类型用
\passthrough{\lstinline!type()!}，查看结果用
\passthrough{\lstinline!print()!}。

\begin{lstlisting}[language=Python]
model_id = "mobileclip_s0"
feature_dim = 512
input_scale = 0.00392156862745098
is_admitted = True

print(type(model_id))
print(type(feature_dim))
print(type(input_scale))
print(type(is_admitted))
\end{lstlisting}

输出：

\begin{lstlisting}
<class 'str'>
<class 'int'>
<class 'float'>
<class 'bool'>
\end{lstlisting}

\passthrough{\lstinline!str!} 是文本，\passthrough{\lstinline!int!}
是整数，\passthrough{\lstinline!float!}
是小数，\passthrough{\lstinline!bool!} 是布尔值。模型 ID 适合存成
\passthrough{\lstinline!str!}，嵌入维度适合存成
\passthrough{\lstinline!int!}。

\subsection{数字运算与单位换算}\label{ux6570ux5b57ux8fd0ux7b97ux4e0eux5355ux4f4dux6362ux7b97}

模型文件大小经常需要在字节和 MiB 之间换算。\passthrough{\lstinline!/!}
是普通除法，\passthrough{\lstinline!//!}
是向下取整，\passthrough{\lstinline!\%!}
是取余数，\passthrough{\lstinline!round()!} 控制显示精度。

\begin{lstlisting}[language=Python]
bytes_size = 131904474
mib_size = bytes_size / 1024 / 1024

print(mib_size)
print(round(mib_size, 2))
print(bytes_size // 1024)
print(bytes_size % 1024)
\end{lstlisting}

输出：

\begin{lstlisting}
125.79369640350342
125.79
128812
986
\end{lstlisting}

\passthrough{\lstinline!round()!} 只改变显示结果，不改变原始变量。

\subsection{列表与元组}\label{ux5217ux8868ux4e0eux5143ux7ec4}

模型输入形状、预处理均值和标准差通常是有序数值。\passthrough{\lstinline!list!}
保存可增删的有序数据，\passthrough{\lstinline!tuple!}
保存创建后不修改的固定组合。

\begin{lstlisting}[language=Python]
input_shape = [1, 3, 224, 224]
image_mean = (0.485, 0.456, 0.406)

print(input_shape[0])
print(input_shape[-1])
print(len(input_shape))
print(image_mean[0])
\end{lstlisting}

输出：

\begin{lstlisting}
1
224
4
0.485
\end{lstlisting}

列表索引从 0 开始，\passthrough{\lstinline!-1!}
表示最后一个元素。\passthrough{\lstinline!len()!} 返回元素个数。

\subsection{字典}\label{ux5b57ux5178}

模型清单由多个字段组成，每个字段用一个键和对应的值表示。\passthrough{\lstinline!dict!}
保存键值对，适合描述一个模型条目。

\begin{lstlisting}[language=Python]
model = {
    "model_id": "mobileclip_s0",
    "input_shape": [1, 3, 256, 256],
    "precision_mode": "mixed_fp16",
    "status": "admitted",
}

print(model["model_id"])
print(model["input_shape"])
\end{lstlisting}

输出：

\begin{lstlisting}
mobileclip_s0
[1, 3, 256, 256]
\end{lstlisting}

键通常使用字符串，值可以是字符串、数字、列表或另一个字典。

\subsection{集合}\label{ux96c6ux5408}

需要知道模型清单中出现过哪些精度模式时，用 \passthrough{\lstinline!set!}
去重。

\begin{lstlisting}[language=Python]
precision_modes = {"mixed_fp16", "allow_fp32_to_fp16", "mixed_fp16"}
print(precision_modes)
\end{lstlisting}

输出：

\begin{lstlisting}
{'allow_fp32_to_fp16', 'mixed_fp16'}
\end{lstlisting}

集合不保留插入顺序，也不能通过索引取值。顺序不重要的去重任务才适合使用集合。

\section{2. 控制流}\label{ux63a7ux5236ux6d41}

\subsection{条件分支}\label{ux6761ux4ef6ux5206ux652f}

\passthrough{\lstinline!if!}
处理一个条件，\passthrough{\lstinline!elif!}
处理``否则如果''，\passthrough{\lstinline!else!}
是兜底。判断从上往下执行，先命中的分支生效。

\begin{lstlisting}[language=Python]
status = "admitted"

if status == "admitted":
    print("可进入 NPU 服务")
elif status == "candidate":
    print("仅用于离线评估")
else:
    print("状态未知")
\end{lstlisting}

输出：

\begin{lstlisting}
可进入 NPU 服务
\end{lstlisting}

比较运算符 \passthrough{\lstinline!==!} 判断是否相等，结果是一个布尔值。

\subsection{for 循环}\label{for-ux5faaux73af}

\passthrough{\lstinline!for!}
循环依次取出列表中的元素，重复执行缩进块。累加前先设置
\passthrough{\lstinline!total = 0!}，这个初始值必须写在循环外。

\begin{lstlisting}[language=Python]
feature_dims = [512, 1024, 2048]

total = 0
for dim in feature_dims:
    total = total + dim

print(total)
print(min(feature_dims))
print(max(feature_dims))
\end{lstlisting}

输出：

\begin{lstlisting}
3584
512
2048
\end{lstlisting}

\subsection{range、enumerate 与 zip}\label{rangeenumerate-ux4e0e-zip}

\passthrough{\lstinline!range(n)!} 生成从 0 到
\passthrough{\lstinline!n - 1!}
的整数序列。\passthrough{\lstinline!enumerate()!}
同时返回序号和元素，\passthrough{\lstinline!zip()!}
把多个序列按位置配对。

\begin{lstlisting}[language=Python]
for index, dim in enumerate([512, 1024, 2048], start=1):
    print(index, dim)
\end{lstlisting}

输出：

\begin{lstlisting}
1 512
2 1024
3 2048
\end{lstlisting}

\begin{lstlisting}[language=Python]
model_ids = ["mobileclip_s0", "resnet50_feature"]
dims = [512, 2048]

for model_id, dim in zip(model_ids, dims):
    print(model_id, dim)
\end{lstlisting}

输出：

\begin{lstlisting}
mobileclip_s0 512
resnet50_feature 2048
\end{lstlisting}

\subsection{while、break 与 continue}\label{whilebreak-ux4e0e-continue}

\passthrough{\lstinline!while!}
在条件仍为真时重复执行。循环内必须修改条件变量，否则会无限循环。\passthrough{\lstinline!continue!}
跳过本次循环，\passthrough{\lstinline!break!} 立即结束整个循环。

\begin{lstlisting}[language=Python]
attempt = 0

while attempt < 3:
    print(f"attempt {attempt + 1}")
    attempt += 1
\end{lstlisting}

输出：

\begin{lstlisting}
attempt 1
attempt 2
attempt 3
\end{lstlisting}

\begin{lstlisting}[language=Python]
dims = [512, 1024, 2048]

for dim in dims:
    if dim == 1024:
        continue
    print(dim)

for dim in dims:
    if dim > 1024:
        break
    print(dim)
\end{lstlisting}

输出：

\begin{lstlisting}
512
2048
512
1024
\end{lstlisting}

\subsection{列表与字典推导式}\label{ux5217ux8868ux4e0eux5b57ux5178ux63a8ux5bfcux5f0f}

推导式用一行代码生成新列表或字典，适合结构简单的过滤和转换。

\begin{lstlisting}[language=Python]
dims = [512, 1024, 2048]
large_dims = [dim for dim in dims if dim >= 1024]

print(large_dims)
\end{lstlisting}

输出：

\begin{lstlisting}
[1024, 2048]
\end{lstlisting}

\begin{lstlisting}[language=Python]
sizes = {"mobileclip_s0": 131904474, "resnet50_feature": 102400000}
sizes_mib = {model: round(size / 1024 / 1024, 2) for model, size in sizes.items()}

print(sizes_mib)
\end{lstlisting}

输出：

\begin{lstlisting}
{'mobileclip_s0': 125.79, 'resnet50_feature': 97.66}
\end{lstlisting}

\section{3.
函数与模型清单}\label{ux51fdux6570ux4e0eux6a21ux578bux6e05ux5355}

\subsection{函数封装}\label{ux51fdux6570ux5c01ux88c5}

函数把重复逻辑打包成可复用单元。下面函数接收模型列表，返回按精度模式统计的字典。

\begin{lstlisting}[language=Python]
def count_by_precision(models):
    result = {}
    for model in models:
        mode = model["precision_mode"]
        result[mode] = result.get(mode, 0) + 1
    return result

models = [
    {"precision_mode": "mixed_fp16"},
    {"precision_mode": "allow_fp32_to_fp16"},
    {"precision_mode": "allow_fp32_to_fp16"},
]

summary = count_by_precision(models)
for mode, count in sorted(summary.items()):
    print(f"{mode}: {count}")
\end{lstlisting}

输出：

\begin{lstlisting}
allow_fp32_to_fp16: 2
mixed_fp16: 1
\end{lstlisting}

\passthrough{\lstinline!dict.get(key, 0)!} 在键不存在时返回默认值
0。\passthrough{\lstinline!sorted()!} 让输出顺序稳定。

\subsection{默认参数}\label{ux9ed8ux8ba4ux53c2ux6570}

默认参数让调用更简洁，同时允许按名字覆盖默认值。

\begin{lstlisting}[language=Python]
def format_model(model_id, precision="mixed_fp16"):
    return f"{model_id}: {precision}"

print(format_model("mobileclip_s0"))
print(format_model("resnet50_feature", precision="allow_fp32_to_fp16"))
\end{lstlisting}

输出：

\begin{lstlisting}
mobileclip_s0: mixed_fp16
resnet50_feature: allow_fp32_to_fp16
\end{lstlisting}

\subsection{*args 与 **kwargs}\label{args-ux4e0e-kwargs}

\passthrough{\lstinline!*values!}
收集位置参数，\passthrough{\lstinline!**options!} 收集关键字参数。

\begin{lstlisting}[language=Python]
def summarize(title, *values, **options):
    print(title, values, options)

summarize("models", 512, 1024, status="admitted")
\end{lstlisting}

输出：

\begin{lstlisting}
models (512, 1024) {'status': 'admitted'}
\end{lstlisting}

\subsection{lambda}\label{lambda}

\passthrough{\lstinline!lambda!} 适合写一个短小的临时函数。

\begin{lstlisting}[language=Python]
models = [
    {"model_id": "mobileclip_s0", "embedding_dim": 512},
    {"model_id": "resnet50_feature", "embedding_dim": 2048},
]

models.sort(key=lambda model: model["embedding_dim"], reverse=True)
print([model["model_id"] for model in models])
\end{lstlisting}

输出：

\begin{lstlisting}
['resnet50_feature', 'mobileclip_s0']
\end{lstlisting}

\subsection{f-string}\label{f-string}

f-string 用 \passthrough{\lstinline!f"..."!} 开头，大括号内插入变量。

\begin{lstlisting}[language=Python]
model_id = "mobileclip_s0"
embedding_dim = 512

print(f"{model_id}: {embedding_dim} dim")
\end{lstlisting}

输出：

\begin{lstlisting}
mobileclip_s0: 512 dim
\end{lstlisting}

\subsection{类型注解}\label{ux7c7bux578bux6ce8ux89e3}

类型注解帮助阅读代码，也方便工具检查错误。它不会改变运行结果。

\begin{lstlisting}[language=Python]
def admitted_models(models: list[dict]) -> list[dict]:
    return [model for model in models if model.get("status") == "admitted"]
\end{lstlisting}

\subsection{dataclass}\label{dataclass}

\passthrough{\lstinline!dataclass!}
自动生成初始化、比较和打印方法，适合保存结构清晰的记录。

\begin{lstlisting}[language=Python]
from dataclasses import dataclass

@dataclass
class ModelRecord:
    model_id: str
    embedding_dim: int
    status: str = "admitted"

record = ModelRecord("mobileclip_s0", 512)
print(record)
\end{lstlisting}

输出：

\begin{lstlisting}
ModelRecord(model_id='mobileclip_s0', embedding_dim=512, status='admitted')
\end{lstlisting}

\subsection{class 与 self}\label{class-ux4e0e-self}

\passthrough{\lstinline!\_\_init\_\_!}
是实例创建时执行的方法，\passthrough{\lstinline!self!} 指向当前实例。

\begin{lstlisting}[language=Python]
class ModelChecker:
    def __init__(self, status="candidate"):
        self.status = status

    def is_admitted(self):
        return self.status == "admitted"

checker = ModelChecker("admitted")
print(checker.is_admitted())
\end{lstlisting}

输出：

\begin{lstlisting}
True
\end{lstlisting}

\section{4.
文件、JSON、命令行与异步}\label{ux6587ux4ef6jsonux547dux4ee4ux884cux4e0eux5f02ux6b65}

\subsection{缺失键与类型转换}\label{ux7f3aux5931ux952eux4e0eux7c7bux578bux8f6cux6362}

模型清单可能来自不同版本的脚本，字段并不总是一致。直接读取缺失字段会触发
\passthrough{\lstinline!KeyError!}，用
\passthrough{\lstinline!dict.get()!} 可以提供默认值。

\begin{lstlisting}[language=Python]
model = {
    "model_id": "mobileclip_s0",
    "precision_mode": "mixed_fp16",
}

print(model.get("status", "unknown"))
\end{lstlisting}

输出：

\begin{lstlisting}
unknown
\end{lstlisting}

把字符串转成数字时，输入内容决定转换是否成功。\passthrough{\lstinline!int("fp16")!}
会触发 \passthrough{\lstinline!ValueError!}，可以用
\passthrough{\lstinline!try/except!} 捕获。

\begin{lstlisting}[language=Python]
value = "fp16"

try:
    numeric_value = int(value)
except ValueError:
    numeric_value = None

print(numeric_value)
\end{lstlisting}

输出：

\begin{lstlisting}
None
\end{lstlisting}

\passthrough{\lstinline!None!}
表示没有有效数值。统计前应先检查它，不要把
\passthrough{\lstinline!None!} 当成 0。

\subsection{pathlib 与 with open}\label{pathlib-ux4e0e-with-open}

\passthrough{\lstinline!Path!}
提供跨平台的路径操作，\passthrough{\lstinline!with open!}
在读取或写入结束后自动关闭文件。

\begin{lstlisting}[language=Python]
from pathlib import Path

path = Path("tmp/appendix3/notes.txt")
path.parent.mkdir(parents=True, exist_ok=True)
path.write_text("mixed_fp16\n", encoding="utf-8")

with path.open(encoding="utf-8") as f:
    print(f.read().strip())
\end{lstlisting}

输出：

\begin{lstlisting}
mixed_fp16
\end{lstlisting}

\subsection{json.loads 与
json.dumps}\label{json.loads-ux4e0e-json.dumps}

模型配置常以 JSON 保存。\passthrough{\lstinline!json.loads()!} 把 JSON
文本转回 Python 对象，\passthrough{\lstinline!json.dumps()!} 把 Python
对象转成 JSON 文本。

\begin{lstlisting}[language=Python]
import json

text = '{"model_id": "mobileclip_s0", "status": "admitted"}'
data = json.loads(text)

print(data["model_id"])
print(data.get("status", "unknown"))
\end{lstlisting}

输出：

\begin{lstlisting}
mobileclip_s0
admitted
\end{lstlisting}

\begin{lstlisting}[language=Python]
config = {"model_id": "mobileclip_s0", "precision_mode": "mixed_fp16"}
print(json.dumps(config, ensure_ascii=False, indent=2))
\end{lstlisting}

输出：

\begin{lstlisting}
{
  "model_id": "mobileclip_s0",
  "precision_mode": "mixed_fp16"
}
\end{lstlisting}

解析失败时，\passthrough{\lstinline!json.loads()!} 会抛出
\passthrough{\lstinline!json.JSONDecodeError!}。处理外部配置或日志时，先确认它确实是合法
JSON，再读取字段，不要把解析失败和字段缺失混为一谈。

\subsection{try / except / else /
finally}\label{try-except-else-finally}

\passthrough{\lstinline!else!}
在没有异常时执行，\passthrough{\lstinline!finally!}
无论是否异常都会执行。

\begin{lstlisting}[language=Python]
try:
    dim = int("512")
except ValueError:
    dim = None
else:
    print("converted")
finally:
    print("done")
\end{lstlisting}

输出：

\begin{lstlisting}
converted
done
\end{lstlisting}

\subsection{argparse}\label{argparse}

\passthrough{\lstinline!argparse!}
用来解析命令行参数，样例脚本中常用来接收模型 ID、精度和路径。

\begin{lstlisting}[language=Python]
import argparse

parser = argparse.ArgumentParser()
parser.add_argument("--model-id", default="mobileclip_s0")
args = parser.parse_args(["--model-id", "resnet50_feature"])

print(args.model_id)
\end{lstlisting}

输出：

\begin{lstlisting}
resnet50_feature
\end{lstlisting}

\subsection{\texorpdfstring{if \textbf{name} ==
``\textbf{main}''}{if name == ``main''}}\label{if-name-main}

被直接运行时会执行
\passthrough{\lstinline!main()!}，被其他模块导入时不会自动执行。

\begin{lstlisting}[language=Python]
def main():
    print("check complete")

if __name__ == "__main__":
    main()
\end{lstlisting}

输出：

\begin{lstlisting}
check complete
\end{lstlisting}

\subsection{logging}\label{logging}

\passthrough{\lstinline!logging!} 比 \passthrough{\lstinline!print()!}
更适合记录运行过程，样例服务中常用来区分普通日志和错误日志。

\begin{lstlisting}[language=Python]
import logging

logging.basicConfig(level=logging.INFO)
logger = logging.getLogger("appendix3")
logger.info("model check complete")
\end{lstlisting}

输出：

\begin{lstlisting}
INFO:appendix3:model check complete
\end{lstlisting}

\subsection{subprocess}\label{subprocess}

\passthrough{\lstinline!subprocess.run()!}
启动外部命令，\passthrough{\lstinline!capture\_output=True!}
捕获标准输出。

\begin{lstlisting}[language=Python]
import subprocess

result = subprocess.run(["python", "--version"], capture_output=True, text=True)
print(result.stdout.strip())
\end{lstlisting}

输出内容取决于本机 Python 版本。

\subsection{async / await}\label{async-await}

\passthrough{\lstinline!async def!}
定义异步函数，\passthrough{\lstinline!await!} 等待异步结果。样例中的
WebRTC 和聊天服务大量使用这种写法。

\begin{lstlisting}[language=Python]
import asyncio

async def check():
    return "admitted"

async def main():
    print(await check())

asyncio.run(main())
\end{lstlisting}

输出：

\begin{lstlisting}
admitted
\end{lstlisting}

\section{5. 调试与验证}\label{ux8c03ux8bd5ux4e0eux9a8cux8bc1}

\subsection{最小调试手段}\label{ux6700ux5c0fux8c03ux8bd5ux624bux6bb5}

出错时先确认对象的真实内容，不要只看报错文字。\passthrough{\lstinline!repr()!}
可以显示字符串中的空格和引号。

\begin{lstlisting}[language=Python]
status = " admitted"

print(status)
print(repr(status))
\end{lstlisting}

输出：

\begin{lstlisting}
 admitted
' admitted'
\end{lstlisting}

如果逻辑复杂，可以在可疑位置使用 \passthrough{\lstinline!breakpoint()!}
进入调试器，或把中间变量打印出来。脚本保存后，先在仓库根目录执行：

\begin{lstlisting}[language=bash]
python -m py_compile 脚本路径.py
\end{lstlisting}

\passthrough{\lstinline!py\_compile!} 只检查语法，不运行板端逻辑。昇腾
310B 的 CANN、ATC 和 ACL 检查仍应在真实开发板上完成。

\subsection{验证清单}\label{src-appendix3-validation}

\begin{itemize}
\tightlist
\item[$\square$]
  能说出
  \passthrough{\lstinline!int!}、\passthrough{\lstinline!float!}、\passthrough{\lstinline!str!}、\passthrough{\lstinline!bool!}
  的区别。
\item[$\square$]
  能用列表保存输入形状或特征维度，并用循环完成求和。
\item[$\square$]
  能用字典表示一个模型条目，按字段名读取值。
\item[$\square$]
  能写函数统计模型清单中的精度模式。
\item[$\square$]
  能区分
  \passthrough{\lstinline!KeyError!}、\passthrough{\lstinline!ValueError!}
  和 \passthrough{\lstinline!json.JSONDecodeError!}。
\item[$\square$]
  能阅读 \passthrough{\lstinline!samples/!}
  中常见的列表推导、\passthrough{\lstinline!enumerate!}、\passthrough{\lstinline!argparse!}
  和 \passthrough{\lstinline!pathlib!} 写法。
\item[$\square$]
  示例只使用标准库，没有 \passthrough{\lstinline!pandas!}
  或板端运行时依赖。
\end{itemize}

\subsection{常见错误与坑}\label{src-appendix3-pitfalls}

{ % do not increment counter
\begin{longtable}[]{@{}
  >{\raggedright\arraybackslash}p{(\linewidth - 4\tabcolsep) * \real{0.3333}}
  >{\raggedright\arraybackslash}p{(\linewidth - 4\tabcolsep) * \real{0.3333}}
  >{\raggedright\arraybackslash}p{(\linewidth - 4\tabcolsep) * \real{0.3333}}@{}}
\toprule\noalign{}
\begin{minipage}[b]{\linewidth}\raggedright
现象
\end{minipage} & \begin{minipage}[b]{\linewidth}\raggedright
可能原因
\end{minipage} & \begin{minipage}[b]{\linewidth}\raggedright
处理方式
\end{minipage} \\
\midrule\noalign{}
\endhead
\bottomrule\noalign{}
\endlastfoot
读取缺失字段报 \passthrough{\lstinline!KeyError!} & 模型清单字段不一致 &
先审计字段，再用 \passthrough{\lstinline!get()!} \\
\passthrough{\lstinline!int("fp16")!} 报错 & 字符串不是数字 & 用
\passthrough{\lstinline!try/except!} 捕获并记录 \\
JSON 解析失败 & 文本不是合法 JSON & 检查完整输入，捕获
\passthrough{\lstinline!JSONDecodeError!} \\
统计结果比预期少 & 把字符串当成数字 & 先看
\passthrough{\lstinline!type()!}，必要时转换 \\
输出顺序每次不同 & 遍历集合或字典 & 需要稳定顺序时用
\passthrough{\lstinline!sorted()!} \\
\end{longtable}
}

\section{6. 作业}\label{src-appendix3-assignment}

\begin{enumerate}
\def\labelenumi{\arabic{enumi}.}
\tightlist
\item
  新建一个脚本，统计
  \passthrough{\lstinline!feature\_dims = [512, 1024, 2048, 512]!}
  的元素个数、最小值、最大值和去重后的维度集合。
\item
  写一个函数
  \passthrough{\lstinline!count\_admitted(models)!}，接收模型字典列表，返回
  \passthrough{\lstinline!status == "admitted"!} 的模型数量。
\item
  写一个 \passthrough{\lstinline!safe\_get(data, key)!}
  函数：键存在时返回值，不存在或类型错误时返回
  \passthrough{\lstinline!None!}。
\item
  用列表推导式筛出 \passthrough{\lstinline![512, 1024, 2048, 512]!}
  中大于等于 1024 的元素。
\item
  用 \passthrough{\lstinline!argparse!} 写一个脚本，接收
  \passthrough{\lstinline!--model-id!}，并输出该参数值。
\item
  把上面的脚本都交给 DSH 或人工审查一次，逐行解释每个分支和异常处理。
\end{enumerate}

AI 辅助审查时仍遵循先理解计划、检查命令、小步执行、验证结果。DSH
生成代码后，必须自己运行并解释输出，不能只看``已完成''的结论。
